\section{Introduction}\label{sec:intro}

Camera-equipped smart glasses are a fast-growing device category~\cite{idc2026xr,counterpoint2026aiglasses}, and an assistant on them is most useful when the wearer can refer to things in view. Pointing is the most direct way to do so: ``what is \emph{that}?'', ``remind me to buy \emph{this}''. Interfaces have combined speech with pointing since \emph{Put-that-there}~\cite{bolt1980put}, and recent wearable assistants use pointing and gaze to resolve such references~\cite{lee2024gazepointar}. All of them first need the answer to a spatial question, from the glasses' own camera: \emph{which object is the wearer pointing at?}

The models that answer this question best today are large. Multimodal LLMs reason from the hand to a ray and then to the target~\cite{egopoint,li2026pointvgr}, and the most accurate grounding model is a transformer of more than 226M parameters with explicit ray modules~\cite{li2026vistaref}. Such models cannot run on glasses, so every query would send camera frames to a phone or a server. That spends what glasses have least of. Their batteries hold a few hundred milliampere-hours, shared by the camera, audio and radio~\cite{yuan2026aiglasses}; every frame sent adds network latency and server cost; and the frames carry images of bystanders off the device. A model that runs on the glasses avoids all of these costs, provided it is accurate enough to be useful. Existing on-device work stops short of this task: egocentric fingertip systems, one of them on a phone, stop at the fingertip~\cite{huang2016pointing,jain2019gestarlite}, and tiny detectors on smart glasses find generic objects, not the one a gesture refers to~\cite{moosmann2024tinyissimoyolo}. Two difficulties are specific to glasses. The camera sees the wearer's hand from behind and often only in part, and other people's hands can be in view and must not be mistaken for the wearer's.

We therefore ask how much of this task a model small enough for the glasses can do. SpatialNet takes one RGB frame and nothing else: no text prompt, no hand box, no depth, and no external hand detector. In one forward pass of a 0.72M-parameter convolutional network, it predicts the wearer's hand, the fingertip, the pointing direction and the box of the pointed-at object, together with a hand-presence gate. Two design choices address the difficulties above. The prediction heads form a chain in which each head reads a shared feature map at the location that the previous one predicted, so that all outputs belong to the same hand (\emph{instance-anchored heads}). And training pastes other people's hands into the frames without labels and penalises any response to them, so that the model learns to follow the wearer's hand without labelled examples of bystanders.

We evaluate SpatialNet on the EgoPoint benchmark~\cite{egopoint} against published large models and against compact detectors that we train at the same compute and score with the same code, and we deploy it on a pre-production pair of smart glasses. We do not claim a more accurate localiser than large models, or tighter boxes than a compact detector. We show what a sub-million-parameter, image-only model can do on glasses, where it falls short, and which design choices matter at this scale.

\paragraph{Contributions.}
\begin{enumerate}
\item \emph{An image-only pointing-target localiser that runs on smart glasses.} A 0.72M-parameter CNN, at least 300 times smaller than published pointing-grounding models, returns the wearer's hand, fingertip, direction and target box in one pass and runs in 60--80\,ms per frame on a pre-production glasses unit (Sections~\ref{sec:method} and~\ref{sec:ondevice}).
\item \emph{A controlled comparison with compact detectors at the same compute.} Detectors draw tighter boxes, but SpatialNet picks the labelled object more often in every controlled comparison, and an off-the-shelf recogniser names the object from its crops as often as from a detector's (Sections~\ref{sec:comparison} and~\ref{sec:choice}).
\item \emph{Evidence on what matters for small pointing models.} Input resolution is the largest lever for accuracy; instance anchoring makes the hand box and the direction more precise; training with distractor hands keeps accuracy when other people's hands are pasted into the frame, and does so for a generic detector too; and an explicit ray prior is not needed at this scale (Section~\ref{sec:design}). We also find that at least 41\% of the EgoPoint test images are COCO photographs, which benchmarks on this dataset should take into account (Section~\ref{sec:data}).
\end{enumerate}
